\documentclass[12pt,a4paper]{article}
\usepackage{ucebnice}

\title{Základní matematické poznatky}
\author{Ondřej Škubala}
\date{\today}

\begin{document}

\maketitle
\newpage



\section{Výroková logika}

\begin{quotebox}{„Všichni lidé jsou smrtelní.\\
Sokrates je člověk.\\
\emph{Tedy Sokrates je smrtelný.“}}{Aristoteles}
\end{quotebox}

\subsection{Výrok}

V běžném životě neustále vyslovujeme nejrůznější věty. Některé z těchto vět mají podobu otázek, jiné rozkazů či přání.  
\important{Matematická logika} se však zaměřuje pouze na takové věty, o nichž můžeme jednoznačně říct, zda jsou pravdivé, nebo nepravdivé.  
Právě tyto věty nazýváme \important{výroky}.

\begin{definitionbox}
\textbf{Výrok} je oznamovací věta, u které má smysl rozhodnout, zda je \important{pravdivá}, nebo \important{nepravdivá}.  
A tuto skutečnost o ni lze jednoznačně prohlásit.
\end{definitionbox}

\begin{examplebox}
\textbf{Co výrok je:}
\begin{itemize}
    \item „Hlavní město Indie je Staré Dillí.“                          \hfill (nepravdivý výrok)
    \item „Číslo $\pi$ je iracionální.“                                 \hfill (pravdivý výrok)
    \item „Pro všechna $a,b$ platí, že $(a + b)^2 = a^2 + 2ab + b^2 $“   \hfill (pravdivý výrok)
\end{itemize}

\textbf{Co výrok není:}
\begin{itemize}
    \item „Listí a strom“                           \hfill (nedává smysl jako tvrzení)
    \item „Kolik je hodin?“                         \hfill (otázka)
    \item „Otevři dveře!“                           \hfill (rozkaz)
    \item $(a + b)^2 = a^2 + 2ab + b^2$             \hfill (rovnice)
\end{itemize}
\end{examplebox}
Výroky tedy nejsou otázky, básnické obrazy, nesmysly apod. Může se zdát zvláštní, že 3. výrok: 
\emph{„Pro všechna $a,b$ platí, že $(a + b)^2 = a^2 + 2ab + b^2 $“}, je výrok a dokonce pravdivý avšak
bezeslov zaspán tentýž výrok $(a + b)^2 = a^2 + 2ab + b^2$, již výrokem není. 
Vycházíme totiž z toho faktu, že nevíme co jsou proměnné $a, b$ a pro které hodnoty tato rovnost nastává.  

\newpage
\subsection{Pravdivostní hodnota}

Každý výrok, o kterém lze říct, zda platí, dostává jednu ze dvou \important{pravdivostních hodnot}.  
Pokud výrok skutečně odpovídá realitě, považujeme jej za \emph{pravdivý}.  
Pokud naopak výrok skutečnosti neodpovídá, nazáváme jej \emph{nepravdivý}. 
Pravdivostní hodnoty zapisujeme zjednodušeně pomocí číslic:
\begin{itemize}
    \item \textbf{1} – výrok je pravdivý,  
    \item \textbf{0} – výrok je nepravdivý.
\end{itemize}
Nebo je také zapisujeme jako \emph{P/Np}, nebo \emph{T/F}, z anglického \emph{True} a \emph{False}.
Avšak v matematice se nejčastěji setkáme se zápisem pomocí číslic.

\begin{examplebox}
\begin{itemize}
    \item „Venku svítí slunce.“ $\;\Rightarrow\;$ může mít hodnotu 1 (pokud svítí), nebo 0 (pokud nesvítí).  
    \item „10 je liché číslo.“ $\;\Rightarrow\;$ hodnota 0, protože výrok je nepravdivý.  
\end{itemize}
\end{examplebox}
\subsection{Negace výroku}

Ke každému výroku existuje vždy i \emph{výrok opačné pravdivostní hodnoty}.  
Tím provádíme tzv. \important{negaci výroku}.  
Negace výroku tedy převrací pravdivost: pokud původní výrok platil, pak jeho negace neplatí, a naopak.  

\begin{definitionbox}
Je-li dán výrok $A$, pak \textbf{negace výroku} $A$ je výrok, který je pravdivý právě tehdy, když $A$ není pravdivý.  
Negaci značíme symbolem $\neg A$, případně $\overline{A}$.
\end{definitionbox}

Protože máme výrok o kterém můžeme říct, zdali je pravdivý, či nepravdivý. 
Můžeme z toho také odvodit, že při negaci negovaného výroku, nám nemůže vzniknout 
nová pravdivostní hodnota než byla hodnota výroku původního. Toto pozorování se nazývá \emph{Zákon vyloučení třetí možnosti}.

\begin{notebox}
\textbf{Zákon vyloučení třetí možnosti:}  
Každý výrok je buď pravdivý, nebo nepravdivý. Neexistuje žádná třetí možnost.  
Jinými slovy, platí tedy:
\begin{center}
    $\neg (\neg A) = A$.
\end{center} 
Negace negace výroku je tedy opět původní výrok.
\end{notebox}

\begin{examplebox}
\textbf{Příklady:}
\begin{itemize}
    \item Výrok: „Venku svítí slunce.“  
          Negace: „Venku \emph{nesvítí} slunce.“  
    \item Výrok: „Číslo 10 je liché.“  
          Negace: „Číslo 10 \emph{není} liché.“  
\end{itemize}
\end{examplebox}


Je užitečné zároveň umět negovat výrok, ve kterém se nachází číselný údaj např. nejvýše, resp. alespoň.
Důležité je však uvědomit si, co nám tyto slova říkají.

\begin{notebox}
    Řekneme-li, že se v množině nachází \textbf{alespoň} $k$ prvků, znamená to, že počet prvků je \textbf{větší nebo roven číslu $k$}.

    Řekneme-li, že se v množině nachází \textbf{nanejvýš} $k$ prvků, znamená to, že počet prvků je \textbf{menší nebo roven číslu $k$}.
\end{notebox}

Negaci výroků tohoto typu dostaneme, když si uvědomíme, že  platí:

Neobsahuje-li množina aspoň $k$ prvků, je počet prvků tedy menší než $k$, neboli je menší nebo roven $k-1$. To tedy znamená, že počet prvků takové množiny je nejvýše $k-1$.

Neobsahuje-li množina nanejvýš $k$ prvků, je počet prvků tedy větší než $k$, neboli je větší nebo roven $k+1$. To tedy znamená, že počet prvků takové množiny je nejméně $k+1$.

Z čehož nám poté vyplývá následující a sice, že:
\begin{center}
\renewcommand{\arraystretch}{1.3}
\begin{tabular}{|C{5cm}|C{2cm}|C{7cm}|}
\hline
\rowcolor{pastelblue!40}
\textbf{Výrok} & \textbf{Význam} & \textbf{Negace} \\
\hline
„Množina má \textbf{alespoň} $k$ prvků.“ 
& $n \geq k$
& „Množina má \textbf{nejvýše} $k-1$ prvků.“ \quad ($n \leq k-1$) \\
\hline
„Množina má \textbf{nanejvýš} $k$ prvků.“ 
& $n \leq k$
& „Množina má \textbf{alespoň} $k+1$ prvků.“ \quad ($n \geq k+1$) \\
\hline
„Množina má \textbf{právě} $k$ prvků.“ 
& $n = k$
& „Množina nemá právě $k$ prvků.“ \quad ($n \neq k$) \\
\hline
\end{tabular}
\end{center}

\begin{notebox}
    Z tvrzení, že \emph{„Množina má $k$ prvků“} nelze zjistit, zdali má nanjevýš, nebo alespoň anebo právě $k$ prvků. 
    Z toho důvodu abychom se vyhnuli těmto nejasnostem, musíme vždy ujasnit údaj o počtu prvků v dané množině.
\end{notebox}
\subsection{Složené výroky a výrokové spojky}
Samotné jednoduché výroky by v praxi neměly tak velký význam.  
Proto jednotlivé výroky spojujeme do tzv. \important{složených výroků}, nazývaných také \emph{formule}, pomocí \emph{výrokových spojek}.
A ty pak dále do \important{složených výrokových forem}
Díky nim dokážeme tvořit složitější tvrzení a zkoumat jejich pravdivost.  

\begin{definitionbox}
\textbf{Složený výrok} vzniká spojením dvou nebo více jednoduchých výroků pomocí \important{výrokových spojek}.  
Pravdivost složeného výroku závisí na pravdivostních hodnotách jednotlivých částí.
\end{definitionbox}

\newpage

Máme několik základních výrokových spojek. Avšak nejdůležitějšími typy, kterými se budeme zabývat jsou
\important{konjunkce}, \important{disjunkce}, \important{implikace}, \important{ekvivalence} a řadí se sem i \important{negace}.
Tyto výrokové spojky čteme jako:

\begin{itemize}
    \item \textbf{Konjunkce} ($A \land B$): „$A$ a zároveň $B$“
    \item \textbf{Disjunkce} ($A \lor B$): „$A$ nebo $B$“ (v logice \emph{inkluzivní})
    \item \textbf{Implikace} ($A \Rightarrow B$): „jestliže $A$, pak $B$“
    \item \textbf{Ekvivalence} ($A \Leftrightarrow B$): „$A$ právě tehdy, když $B$“
    \item \textbf{Negace} ($\neg A$): „Není pravda, že …“
\end{itemize}

\subsubsection*{Konjunkce}
\begin{definitionbox}
    \textbf{Konjunkce} libovolných dvou výroků $A, B$ je výrok, který vznikne jejich spojením spojkou \textbf{a}.
    Konjunkci výroků $A, B$ zapisujeme jako $A \land B$ a tento zápis čteme jako „$A$ a $B$“, nebo také jako „$A$ a zároveň $B$“.
\end{definitionbox}


\begin{examplebox}
\begin{tabbing}
    Mějme výroky a z nich utvoříme pomocí konjunkce složené výroky:\\
    $A$: Číslo 5 je prvočíslo. \=
    $B$: Číslo 5 je liché.\\
    $C$: Číslo 5 je záporné.\>
    $D$: Číslo 5 je sudé.\\
    \\
    $A \land B$: Číslo 5 je prvočíslo a zároveň je liché. \\
    $A \land C$: Číslo 5 je prvočíslo a zároveň je záporné.\\
    $D \land B$: Číslo 5 je sudé a zároveň liché.\\
    $D \land C$: Číslo 5 je sudé a zároveň záporné.
\end{tabbing}
\end{examplebox}

Přechozí příklad nám poslouží k tomu, abychom si na něm ukázali jak závisí pravdivostní hodnoty konjunkce na hodnotách jednotlivých výroků.
Ze všech výše vytvořených konjunkcí lze jen o jediné tvrdit, že je pravdivá. A je to hned ta první $A \land B$. O všech ostatních nemůžeme logickou dedukcí
dojít k závěru, že by byli pravdivé. Můžeme tedy říct, že ze všech těchto konjunkcí byla pradivá pouza ta, ve které byly oba výroky pravdivé.
V souladu s touto logickou dedukcí lze nyní definovat pravdivostní hodnotu konjunkce každých dvou výroků takto:

\begin{definitionbox}
    \textbf{Konjunkce} libovolných dvou výroků $A$ a $B$ \textbf{je pravdivá} právě tehdy, jsou-li pradivé oba její výroky $A,B$.
\end{definitionbox}


\subsubsection*{Disjunkce}
\begin{definitionbox}
    \textbf{Disjunkce} libovolných dvou výroků $A, B$ je výrok, který vznikne jejich spojením spojkou \textbf{nebo}.  
    Disjunkci výroků $A, B$ zapisujeme jako $A \lor B$ a tento zápis čteme jako „$A$ nebo $B$“, případně „$A$ alespoň jedno z $A, B$“.
\end{definitionbox}

\begin{notebox}
V matematické logice se výrok „$A$ nebo $B$“ chápe jako \textbf{inkluzivní disjunkce}, tj. pravdivý je i v případě, že platí oba výroky současně.  
Pozor: v běžném jazyce se často „nebo“ používá exkluzivně (buď jedno, nebo druhé, ale ne oboje).
\end{notebox}

\begin{examplebox}
\begin{tabbing}
    Mějme výroky a z nich utvoříme pomocí disjunkce složené výroky:\\
    $A$: Číslo 5 je prvočíslo. \= 
    $B$: Číslo 5 je liché.\\
    $C$: Číslo 5 je záporné.\>
    $D$: Číslo 5 je sudé.\\
    \\
    $A \lor B$: Číslo 5 je prvočíslo nebo je liché. \\ 
    $A \lor C$: Číslo 5 je prvočíslo nebo je záporné.\\
    $D \lor B$: Číslo 5 je sudé nebo je liché.\\
    $D \lor C$: Číslo 5 je sudé nebo je záporné.
\end{tabbing}
\end{examplebox}

Z uvedených příkladů si můžeme všimnout, že disjunkce je pravdivá tehdy, pokud platí alespoň jeden z výroků.  
Například disjunkce $A \lor B$ je pravdivá, protože oba výroky $A$ i $B$ jsou pravdivé. Složený výrok $A \lor C$ je taktéž pravdivý, protože i když $C$ je nepravdivý výrok, tak výrok $A$ je pravdivý.  
Naopak disjunkce $D \lor C$ není pravdivá, protože oba výroky $D$ i $C$ jsou nepravdivé.

\begin{definitionbox}
    \textbf{Disjunkce} libovolných dvou výroků $A$ a $B$ \textbf{je pravdivá} právě tehdy, je-li pravdivý alespoň jeden z výroků $A$ nebo $B$.  
    Nepravdivá je pouze v případě, kdy jsou oba výroky $A, B$ nepravdivé.
\end{definitionbox}

\subsubsection*{Implikace}
\begin{definitionbox}
    \textbf{Implikace} libovolných dvou výroků $A, B$ je výrok, který vznikne spojením „jestliže … pak …“.  
    Implikaci zapisujeme jako $A \Rightarrow B$ a čteme jako „jestliže $A$, pak $B$“.
\end{definitionbox}
\newpage
\begin{examplebox}
\begin{tabbing}
    Mějme výroky a z nich utvoříme pomocí implikace složené výroky:\\
    $A$: Číslo 5 je prvočíslo. \= 
    $B$: Číslo 5 je liché.\\
    $C$: Číslo 5 je záporné.\>
    $D$: Číslo 5 je sudé.\\
    \\
    $A \Rightarrow B$: Jestliže je číslo 5 prvočíslo, pak je liché. \\ 
    $A \Rightarrow D$: Jestliže je číslo 5 prvočíslo, pak je sudé. \\
    $D \Rightarrow A$: Jestliže je číslo 5 sudé, pak je prvočíslo.\\
    $C \Rightarrow D$: Jestliže je číslo 5 záporné, pak je sudé.
\end{tabbing}
\end{examplebox}

Na těchto příkladech odhalíme, že všechny implikace jsou pravdivé, kromě situace $A \Rightarrow D$, která je nepravdivá, protože 5 je sice prvočíslo, ale není sudé.  

\begin{definitionbox}
    \textbf{Implikace} $A \Rightarrow B$ je \textbf{nepravdivá} pouze v případě, že výrok $A$ je pravdivý a výrok $B$ je nepravdivý.  
    Ve všech ostatních případech je pravdivá.
\end{definitionbox}

\begin{notebox}
    U implikace $A \Rightarrow B$ se výrok $A$ nazývá \textbf{předpoklad} (nebo \textbf{antecedent})  
    a výrok $B$ se nazývá \textbf{závěr} (nebo \textbf{konsekvent}).  

    Místo slovního spojení „jestliže $A$, pak $B$“ se také říká, že:  
    \begin{itemize}
        \item $A$ je \textbf{postačující podmínka} pro $B$,  
        \item $B$ je \textbf{nutná podmínka} pro $A$.
    \end{itemize}
\end{notebox}

    Někdy nás zajímá i tzv. \textbf{obrácená implikace}.  
    Pokud platí $A \Rightarrow B$, neznamená to automaticky, že platí i $B \Rightarrow A$.  
    Pokud je výrok $A$ pravdivý a výrok $B$ nepravdivý, pak implikace $A \Rightarrow B$ neplatí.  
    Obrácená implikace $B \Rightarrow A$ však může mít jinou pravdivostní hodnotu, a proto je potřeba je rozlišovat.

\begin{examplebox}
    Mějme dva výroky, výrok $A$: „Trojúhelník je rovnostranný.“ a výrok $B$: „Trojúhelník je rovnoramenný.“

\begin{itemize}
    \item $A \Rightarrow B$: „Jestliže je trojúhelník rovnostranný, pak je rovnoramenný.“  
          \hfill (pravdivý výrok, protože každý rovnostranný trojúhelník je i rovnoramenný)
    \item $B \Rightarrow A$: „Jestliže je trojúhelník rovnoramenný, pak je rovnostranný.“  
          \hfill (nepravdivý výrok, protože existují rovnoramenné trojúhelníky, které nejsou rovnostranné)
\end{itemize}

Tento příklad jasně ukazuje, že \textbf{implikace a její obrácený tvar nejsou totéž}.  
První implikace ($A \Rightarrow B$) platí, ale její obrácená ($B \Rightarrow A$) už nikoli.
\end{examplebox}

Logické spojky můžeme použít i na výroky, které spolu na první pohled nijak nesouvisejí.  
    Logika se totiž zajímá jen o jejich pravdivostní hodnoty, nikoli o obsahovou souvislost. Jako můžeme vidět na příkladu další straně. \newpage

\begin{examplebox}
       Mějme dva pravdivé výroky $A$: „Velryba je savec.“ a $B$: „J. A. Komenský se narodil roku 1592.“.


    \begin{itemize}
        \item $A \Rightarrow B$: „Je-li je velryba savec, pak se J. A. Komenský narodil roku 1592.“  
        \item $B \Rightarrow A$: „Jestliže se J. A. Komenský narodil roku 1592, pak je velryba savec.“
    \end{itemize}

    Přestože spojení působí nesmyslně, z hlediska výrokové logiky se s ním pracuje zcela stejně jako s jinými implikacemi.
    I přesto, že obě implikace jsou pravdivé, tak se naštěstí s podobnými výroky nesetkáváme.
\end{examplebox}

\subsubsection*{Tautologie a kontradikce}

Některé výrokové formule mají tu vlastnost, že jsou \textbf{vždy pravdivé}, bez ohledu na pravdivostní hodnotu jednotlivých výroků.  
Takovým formulím říkáme \important{tautologie}.  

Naopak výrokové formule, které jsou \textbf{vždy nepravdivé}, se nazývají \important{kontradikce}.  

\begin{examplebox}
\begin{itemize}
    \item Tautologie: $A \Rightarrow (A \lor B)$  
          Tento výrok je vždy pravdivý, protože pokud $A$ platí, pak $A \lor B$ je také pravdivé,  
          a pokud $A$ neplatí, pak implikace $A \Rightarrow (A \lor B)$ je stejně pravdivá.  

    \item Kontradikce: $A \land \neg A$  
          Tento výrok je vždy nepravdivý, protože výrok nemůže být současně pravdivý i nepravdivý.
\end{itemize}
\end{examplebox}


\subsubsection*{Ekvivalence}
Nejprve si připomeňme tabulky pro základní spojky:

\begin{center}
\renewcommand{\arraystretch}{1.2}
\begin{tabular}{|C{1cm}|C{1cm}|C{1.5cm}|C{1.5cm}|C{1.5cm}|}
\rowcolor{pastelblue!40}
\hline
$A$ & $B$ & $A \land B$ & $A \lor B$ & $A \Rightarrow B$ \\
\hline
1 & 1 & 1 & 1 & 1 \\
1 & 0 & 0 & 1 & 0 \\
0 & 1 & 0 & 1 & 1 \\
0 & 0 & 0 & 0 & 1 \\
\hline
\end{tabular}
\end{center}

Pokud vezmeme současně $A \Rightarrow B$ i $B \Rightarrow A$, můžeme je spojit do jednoho složeného výroku pomocí konjunkce, který poté nazýváme ekvivalence:  
\[
(A \Rightarrow B) \land (B \Rightarrow A)
\]

\begin{center}
\renewcommand{\arraystretch}{1.2}
\begin{tabular}{|C{1cm}|C{1cm}|C{1.5cm}|C{1.5cm}|C{4cm}|}
\rowcolor{pastelblue!40}
\hline
$A$ & $B$ & $A \Rightarrow B$ & $B \Rightarrow A$ & $(A \Rightarrow B) \land (B \Rightarrow A)$ \\
\hline
1 & 1 & 1 & 1 & 1\\
1 & 0 & 0 & 1 & 0\\
0 & 1 & 1 & 0 & 0\\
0 & 0 & 1 & 1 & 1\\
\hline
\end{tabular}
\end{center}

\begin{definitionbox}
\textbf{Ekvivalence} dvou výroků $A$ a $B$ zapisujeme jako $A \Leftrightarrow B$.  
Čteme „$A$ právě tehdy, když $B$“.
\end{definitionbox}

\begin{notebox}
Někdy se také říká, že „$A$ je nutná a postačující podmínka pro $B$“.  
Toto vyjádření se používá hlavně v matematických definicích a důkazech.
\end{notebox}

Z předchozí pravdivostní tabulky můžeme vidět, že: 

\begin{definitionbox}
    Výrok $A \Leftrightarrow B$ je \textbf{pravdivý} tehdy, když oba výroky $A$ a $B$ mají stejnou pravdivostní hodnotu  
    (oba jsou pravdivé, nebo oba jsou nepravdivé).  
    Nepravdivý je tehdy, když mají různé pravdivostní hodnoty.
\end{definitionbox}

\begin{examplebox}
    Příkladem pravdivé ekvivalence je výrok:  

    \emph{„Druhá mocnina délky nejdelší strany trojúhelníku $ABC$ je rovna součtu druhých mocnin délek jeho zbývajících stran
    právě tehdy, když je trojúhelník $ABC$ pravoúhlý.“}


    Tento výrok je ekvivalence, protože oba dílčí výroky mají vždy stejnou pravdivostní hodnotu:  
    \begin{itemize}
        \item jsou \textbf{pravdivé}, pokud je trojúhelník $ABC$ pravoúhlý,  
        \item jsou \textbf{nepravdivé}, pokud trojúhelník $ABC$ pravoúhlý není.  
    \end{itemize}
\end{examplebox}

\subsubsection*{Negace složených výroků}

Negace složených výroků se dá dobře pochopit na příkladu.  
Podívejme se na výroky $A$ a $B$ a jejich konjunkci $A \land B$.  
Zkusme zjistit, jak vypadá negace $\neg(A \land B)$ pomocí pravdivostní tabulky:

\begin{center}
\renewcommand{\arraystretch}{1.2}
\begin{tabular}{|C{1cm}|C{1cm}|C{2cm}|C{2cm}|}
\rowcolor{pastelblue!40}
\hline
$A$ & $B$ & $A \land B$ & $\neg(A \land B)$ \\
\hline
1 & 1 & 1 & 0 \\
1 & 0 & 0 & 1 \\
0 & 1 & 0 & 1 \\
0 & 0 & 0 & 1 \\
\hline
\end{tabular}
\end{center}

Vidíme, že ve sloupci $\neg(A \land B)$ máme tři jedničky (ve všech případech kromě situace, kdy jsou oba výroky pravdivé).  
Když tuto pravdivostní hodnotu porovnáme s výroky $\neg A$ a $\neg B$, dostaneme:

\begin{center}
\renewcommand{\arraystretch}{1.2}
\begin{tabular}{|C{1cm}|C{1cm}|C{1cm}|C{1cm}|C{3cm}|}
\rowcolor{pastelblue!40}
\hline
$A$ & $B$ & $\neg A$ & $\neg B$ & $(\neg A \lor \neg B)$ \\
\hline
1 & 1 & 0 & 0 & 0 \\
1 & 0 & 0 & 1 & 1 \\
0 & 1 & 1 & 0 & 1 \\
0 & 0 & 1 & 1 & 1 \\
\hline
\end{tabular}
\end{center}

A zjistíme tedy, že poslední sloupce obou tabulek jsou stejné, což znamená, že platí následující logická ekvivalence:
\[
\neg(A \land B) \Leftrightarrow (\neg A \lor \neg B).
\]
Takto jsme odvodili první z \important{De Morganových zákonů}.  
Podobným způsobem lze dokázat i ostatní:
\begin{notebox}
\textbf{De Morganovy zákony:}
\begin{align*}
\neg(A \lor B) &\Leftrightarrow (\neg A) \land (\neg B), \\
\neg(A \land B) &\Leftrightarrow (\neg A \lor \neg B), \\
\neg(A \Rightarrow B) &\Leftrightarrow A \land \neg B, \\
\neg(A \Leftrightarrow B) &\Leftrightarrow (A \land \neg B) \lor (\neg A \land B).
\end{align*}
\end{notebox}


\subsubsection*{Shrnutí pravdivostních hodnot}

Všechny výrokové spojky, které jsme dosud zavedli, můžeme shrnout do jediné \emph{pravdivostní tabulky}.  
Ta ukazuje, jakou pravdivostní hodnotu má složený výrok pro všechny možné kombinace hodnot výroků $A$ a $B$.

\begin{center}
\renewcommand{\arraystretch}{1.2}
\begin{tabular}{|C{1cm}|C{1cm}|C{1.5cm}|C{1.5cm}|C{1.5cm}|C{1.5cm}|C{1.5cm}|}
\rowcolor{pastelblue!40}
\hline
$A$ & $B$ & $\neg A$ & $A \land B$ & $A \lor B$ & $A \Rightarrow B$ & $A \Leftrightarrow B$ \\
\hline
1 & 1 & 0 & 1 & 1 & 1 & 1 \\
1 & 0 & 0 & 0 & 1 & 0 & 0 \\
0 & 1 & 1 & 0 & 1 & 1 & 0 \\
0 & 0 & 1 & 0 & 0 & 1 & 1 \\
\hline
\end{tabular}
\end{center}

Tato tabulka je základním nástrojem výrokové logiky, protože přehledně shrnuje chování všech běžně používaných spojek.
\subsection{Kvantifikované výroky a jejich negace}

Doposud jsme se zabývali jednotlivými výroky (např. „číslo 5 je liché“).  
V matematice se však často setkáváme s výroky, které se vztahují k celé množině objektů, např. \emph{„Každé přirozené číslo je větší než nula.“}.  
Takovým výrokům říkáme \important{kvantifikované výroky}.  
Kvantifikátory vyjadřují, zda vlastnost platí pro všechny prvky množiny, nebo pouze pro některé prvky.  

Rozlišujeme dva kvantifikátory:  
\important{Univerzální} (velký) kvantifikátor a \important{Existenční} (malý) kvantifikátor.

\begin{definitionbox}
\begin{itemize}
    \item \textbf{Univerzální kvantifikace} ($\forall$): říká, že výrok platí pro všechny prvky dané množiny.  
    Například: „\textbf{Všechna} přirozená čísla jsou větší než $0$. “  
    Zapisujeme symbolem $\forall$.
    \item \textbf{Existenční kvantifikace} ($\exists$): říká, že výrok platí alespoň pro jeden prvek množiny.  
    Například: „\textbf{Existuje alespoň jedno} přirozené číslo, které je větší než $10$. “  
    Zapisujeme symbolem $\exists$.
\end{itemize}
\end{definitionbox}

\begin{examplebox}
\begin{itemize}
    \item $\forall n \in \mathbb{N}: n+1 > n$ \hfill (každé přirozené číslo je menší než jeho následník)  
    \item $\exists n \in \mathbb{N}: n^2 = 16$ \hfill (existuje přirozené číslo, jehož druhá mocnina je 16)
\end{itemize}
\end{examplebox}

Stejně jako ve výrokové logice, i u kvantifikátorů můžeme tvořit \important{negace}.  
Platí jednoduchý princip: negací univerzálního výroku je výrok existenční a negací existenčního výroku je výrok univerzální.

\begin{notebox}
\textbf{Pravidla pro negaci kvantifikovaných výroků:}
\begin{align*}
\neg(\forall x \in M: P(x)) &\Leftrightarrow \exists x \in M: \neg P(x) \\
\neg(\exists x \in M: P(x)) &\Leftrightarrow \forall x \in M: \neg P(x)
\end{align*}
\end{notebox}

Pro přehlednost lze pravidla shrnout do tabulky:
\begin{center}
\renewcommand{\arraystretch}{1.3}
\begin{tabular}{|C{6,5cm}|C{8cm}|}
\hline
\rowcolor{pastelblue!40}
\textbf{Výrok} & \textbf{Negace výroku} \\
\hline
„Každý … je …“ & „Existuje alespoň jeden …, který není …“ \\
\hline
„Alespoň jeden … je …“ & „Pro každý … platí, že není …“ \\
\hline
„Žádný … není …“ & „Existuje alespoň jeden …, který je …“ \\
\hline
\end{tabular}
\end{center}

\begin{examplebox}
\begin{itemize}
    \item Výrok: „Každý žák složil zkoušku.“  
          Negace: „Existuje alespoň jeden žák, který nesložil zkoušku.“  
    \item Výrok: „Existuje reálné číslo, které je větší než 10.“  
          Negace: „Žádné reálné číslo není větší než 10.“  
\end{itemize}
\end{examplebox}

Podobně se negují i výroky s nerovnostmi: 

\begin{notebox}
\begin{align*}
\neg(x > a) &\Leftrightarrow x \leq a \\
\neg(x \geq a) &\Leftrightarrow x < a \\
\neg(x < a) &\Leftrightarrow x \geq a \\
\neg(x \leq a) &\Leftrightarrow x > a \\
\neg(x = a) &\Leftrightarrow x \neq a
\end{align*}
\end{notebox}

\subsection{Logická výstavba matematiky}

Základní výstavba matematiky je založena na \important{axiomech} a jejich soustavách.  
\textbf{Axiomy} jsou základní tvrzení (výroky), které přijímáme jako pravdivé i bez důkazu.  
Axiomy tvoří výchozí bod, z něhož se pak logickými postupy odvozují další matematická tvrzení.  

Axiomy mají několik základních vlastností:
\begin{itemize}
    \item \textbf{Jednoduchost}: axiom by měl být co nejjednodušší a intuitivně pochopitelný,  
    \item \textbf{Bezespornost}: axiom by neměl vést k rozporům,  
    \item \textbf{Nezávislost}: axiom nesmí být odvoditelný z jiných axiomů,  
    \item \textbf{Úplnost}: soustava axiomů by měla umožnit rozhodnout pravdivost nebo nepravdivost dalších tvrzení.
\end{itemize}

\begin{examplebox}
Příklady axiomů:
\begin{itemize}
    \item \textbf{Peanovy axiomy}, které definují přirozená čísla $\mathbb{N}$.  
    \item Axiomy na začátku \textit{Eukleidových Základů}, které tvoří základ klasické geometrie.
\end{itemize}
\end{examplebox}

\begin{definitionbox}
\textbf{Definice} je přesné určení významu pojmů na základě již známých pojmů, pravidel a jejich vlastností.  
Definice musí být jednoznačná a neměla by obsahovat tzv. „logické kruhy“ (definice kruhem).  
\end{definitionbox}

\begin{examplebox}
Příklad definice:  
„Číslo je \textbf{sudé}, pokud existuje celé číslo $k$, pro které platí $n = 2k$.“
\end{examplebox}

Z axiomů a definic se pak odvozují \important{matematické věty}.  
Matematická věta je výrok, který lze dokázat na základě axiomů, definic nebo jiných vět a má konkrétní matematický obsah.  
Matematické věty mají často tvar implikace nebo ekvivalence.  

Rozlišujeme několik typů vět:
\begin{itemize}
    \item \textbf{Lemma}: pomocná věta, která spojuje a pomáhá při důkazech složitějších vět.  
    \item \textbf{Věta}: obecné tvrzení.  
    \item \textbf{Koro­lár (důsledek)}: věta, která přímo vyplývá z jiné věty.
\end{itemize}

\begin{examplebox}
Příklad matematické věty:  
\textbf{Pythagorova věta} – V každém pravoúhlém trojúhelníku platí:  
„Obsah čtverce nad přeponou je roven součtu obsahů čtverců nad odvěsnami.“  
Symbolicky: \[ c^2 = a^2 + b^2. \]
\end{examplebox}

\subsubsection*{Hierarchie matematických tvrzení}

Celá matematika se dá chápat jako logická stavba, kde každý prvek navazuje na ten předchozí:

\begin{center}
\renewcommand{\arraystretch}{1.4}
\begin{tabular}{|C{3cm}|C{11,5cm}|}
\hline
\rowcolor{pastelblue!40}
\textbf{Úroveň} & \textbf{Popis} \\
\hline
\textbf{Axiomy} & Základní tvrzení přijímaná bez důkazu (např. „Pro každé číslo $n$ existuje jeho následník $n+1$“ – Peanův axiom). \\
\hline
\textbf{Definice} & Určení významu pojmů (např. „Sudé číslo je číslo dělitelné dvěma“). \\
\hline
\textbf{Věty} & Tvrzení dokazovaná na základě axiomů a definic (např. Pythagorova věta). \\
\hline
\textbf{Důsledky (koroláry)} & Věty přímo vyplývající z jiných vět (např. z Pythagorovy věty plyne: „Pravoúhlý trojúhelník nemůže být rovnoramenný se stranami délky 3, 3, 3“). \\
\hline
\end{tabular}
\end{center}

\begin{examplebox}
\textbf{Ukázka logické výstavby v analýze:}

\begin{itemize}
    \item \textbf{Axiom}: Každá shora omezená množina reálných čísel má supremum.  
    \item \textbf{Definice}: Posloupnost $(a_n)$ má limitu $A$, jestliže pro každé $\varepsilon > 0$ existuje $N \in \mathbb{N}$ takové, že pro všechna $n \geq N$ platí $|a_n - A| < \varepsilon$.  
    \item \textbf{Lemma}: Má-li posloupnost limitu, pak je tato limita jednoznačná.  
    \item \textbf{Věta}: Pokud $(a_n) \to A$ a $(b_n) \to B$, pak platí $(a_n + b_n) \to A + B$.  
    \item \textbf{Důsledek}: Pokud $(a_n) \to A$ a $c \in \mathbb{R}$, pak $(c \cdot a_n) \to cA$.
\end{itemize}
\end{examplebox}

\subsubsection*{Základní typy matematických důkazů}

Matematické věty dokazujeme různými způsoby. Mezi nejčastější patří:

\begin{definitionbox}
\textbf{Přímý důkaz} spočívá v tom, že z předpokladů postupně odvozujeme další tvrzení, až dojdeme k závěru.  
Typické je dokazování implikace $A \Rightarrow B$: předpokládáme, že $A$ platí, a logickými kroky ukážeme, že z toho nutně plyne i $B$.
\end{definitionbox}

\begin{examplebox}
Příklad přímého důkazu:  
„Jestliže je číslo $n$ sudé, pak je $n^2$ sudé.“  

Předpokládejme, že $n = 2k$ pro nějaké celé $k$. Pak $n^2 = (2k)^2 = 4k^2 = 2 \cdot (2k^2)$, což je opět sudé číslo.  
Tím jsme dokázali implikaci.
\end{examplebox}

\begin{definitionbox}
\textbf{Důkaz sporem} probíhá tak, že předpokládáme negaci tvrzení, které chceme dokázat, a ukážeme, že tento předpoklad vede k logickému rozporu.  
Tím pádem původní tvrzení musí platit.
\end{definitionbox}

\begin{examplebox}
Příklad důkazu sporem:  
„Neexistuje celé číslo $n$, pro které $2n+1$ je sudé.“  

Důkaz sporem: Předpokládejme opak, tj. že existuje celé číslo $n$, pro které je $2n+1$ sudé.  
Pak ale $2n$ je sudé a po přičtení $1$ získáme liché číslo. To je rozpor.  
Tedy $2n+1$ nemůže být sudé.
\end{examplebox}


% --- PŘÍKLADY ---

\exercisepart

\begin{exercise}[Určete, co je a není výrok]
Rozhodněte, zda se jedná o výrok. Pokud ano, určete, zda je jeho pravdivostní hodnota zjistitelná.
\begin{tasks}(1)
\task „Kolik je hodin?“  
\task „Číslo 12 je dělitelné třemi.“  
\task „V roce 2100 bude žít na Zemi více než 12 miliard lidí.“  
\task „Praha je hlavní město České republiky.“  
\task „Každý jednorožec má křídla.“  
\end{tasks}
\end{exercise}

\begin{exercise}[Posouzení pravdivosti výroků]
Určete, zda jsou následující výroky pravdivé nebo nepravdivé:
\begin{tasks}(1)
\task „Číslo $5$ je nezáporné.“  
\task „Odmocnina ze $49$ není $7$.“  
\task „Číslo $6-9$ není kladné.“  
\task „Není nepravda, že pravdivý výrok je vždy nepravdivý.“  
\end{tasks}
\end{exercise}

\begin{exercise}[Negace výroků – jednoduché]
Znegujte následující výroky:
\begin{tasks}(1)
\task „Trojúhelník $ABC$ je pravoúhlý.“  
\task „Trojúhelník $ABC$ je ostroúhlý.“  
\task „Venku prší.“  
\task „$\sqrt{2} + \sqrt{3} \leq \sqrt{5}$.“  
\task „Každé prvočíslo je liché.“  
\task „Sterakdary je nejlepší YouTuber.“
\end{tasks}
\end{exercise}

\begin{exercise}[Negace výroků s kvantifikátory]
Znegujte následující výroky:
\begin{tasks}(1)
\task „Množina $M$ má právě $5$ prvků.“  
\task „Množina $M$ má nanejvýš $10$ prvků.“  
\task „Množina $M$ má alespoň $2$ prvky.“  
\task „V každém trojúhelníku je alespoň jeden úhel ostrý.“  
\task „Existuje přirozené číslo, které je menší než nula.“  
\end{tasks}
\end{exercise}

\newpage

\begin{exercise}[Výroky o kružnicích]
Mějme dvě kružnice $k_1$ a $k_2$. Jsou dány následující výroky:
\begin{tasks}(1)
\task „Kružnice nemají žádný společný bod.“  
\task „Kružnice se protínají.“  
\task „Kružnice mají vnitřní dotek.“  
\end{tasks}

Najděte negaci výroku:  
\begin{center}
„Kružnice mají vnější dotyk.“
\end{center}
\end{exercise}

\begin{exercise}[Konjunkce čtyř výroků]
Mějme čtyři výroky:
\[
\begin{aligned}
a &: \text{„Slunce vychází na východě.“}\\
b &: \text{„Všichni tučňáci umí létat.“}\\
c &: \text{„Praha je hlavní město České republiky.“}\\
d &: \text{„Londýn leží v Austrálii.“}
\end{aligned}
\]

Vytvořte z nich několik konjunkcí a určete jejich pravdivostní hodnotu (stručně zdůvodněte).
\begin{tasks}(4)
\task $a \land c$,
\task $a \land b$,
\task $(a \land c) \land d$,
\task $b \land d$.
\end{tasks}
\end{exercise}


\begin{exercise}[Disjunkce čtyř výroků]
Mějme tytéž výroky z úlohy 1.6. $a, b, c, d$. 
Vytvořte z nich několik disjunkcí a určete jejich pravdivostní hodnotu (stručně zdůvodněte).
U každé disjunkce určete pravdivostní hodnotu a vysvětlete proč (nezapomeňte, že „nebo“ je inkluzivní).
\begin{tasks}(4)
\task $a\lor b$, 
\task $a\lor c$, 
\task $(a\lor c) \lor d$,
\task $(b\lor d)$.
\end{tasks}
\end{exercise}

\begin{exercise}[Vyšetřete pravdivost výroku]
Určete (pomocí pravdivostní tabulky), kdy je pravdivý složený výrok
\[
(\neg A \lor B) \land A.
\]
\begin{tasks}(1)
\task Sestavte pravdivostní tabulku pro všechny kombinace hodnot $A,B\in\{0,1\}$.
\task Vyznačte řádky, ve kterých je celý výraz pravdivý.
\end{tasks}
\end{exercise}

\begin{exercise}[Jazyk vs. logika – „a“ a „nebo“]
Zdůvodněte, proč je výrok \emph{„Praha a Washington jsou evropská města“} nepravdivý, 
zatímco výrok \emph{„Praha nebo Washington jsou evropská města“} je pravdivý.
\begin{tasks}(1)
\task Vysvětlete, proč konjunkce selhává.
\task Vysvětlete, proč je disjunkce pravdivá.
\end{tasks}
\end{exercise}

\begin{exercise}[Vyhodnocení složených výrazů pro dané pravdivosti]
Nechť $A$ je pravdivý výrok, $B$ nepravdivý výrok a $C$ pravdivý výrok. Určete pravdivost následujících složených výroků:
\begin{tasks}(1)
\task $(A \lor B) \land C$
\task $(A \land B) \lor C$
\task $(A \land C) \lor B$
\task $(A \lor C) \land B$
\end{tasks}
\end{exercise}

\begin{exercise}[Závislost pravdivosti na $A$ a $B$]
Prozkoumejte, na jakých hodnotách pro $A,B$ jsou pravdivé následující výroky:
\begin{tasks}(1)
\task $(\neg A \land B) \lor A$
\task $(\neg A \land B) \lor B$
\task $(\neg A \lor B) \lor B$
\end{tasks}

\medskip
Doporučení: Sestavte pro každý výraz pravdivostní tabulku a určete, zda nejde o tautologii (vždy pravda) nebo zda se výraz zjednoduší
(např. pomocí absorbčních a idempotentních zákonů).
\end{exercise}

\begin{exercise}[Implikace a pravdivostní tabulky]
Vyšetřete pomocí pravdivostních tabulek, kdy jsou pravdivé následující implikace:
\begin{tasks}(1)
\task $A \Rightarrow (A \lor B)$
\task $(A \lor B) \Rightarrow A$
\task $(A \lor B) \Rightarrow (A \land B)$
\end{tasks}
\end{exercise}

\begin{exercise}[„Paradoxní“ implikace s reálným tvrzením]
Adam, který není náčelníkem kmene Apačů, prohlásil:  
\emph{„Je-li odmocnina z 10 menší než tři, pak jsem náčelníkem kmene Apačů.“}  
Posuďte pravdivost:
\begin{tasks}(1)
\task původní implikace,
\task \textbf{obrácené} implikace,
\task \textbf{obměněné} (kontrapozice) původní implikace.
\end{tasks}
\end{exercise}

\begin{exercise}[Pravoúhlý úhel a Thaletova kružnice]
Jsou dány dva výroky:  
$a$: „Úhel $ACB$ je pravý.“ \quad $b$: „Bod $C$ leží na kružnici o průměru $AB$.“
Posuďte pravdivost následujících výroků:
\begin{tasks}(1)
\task $a \Rightarrow b$
\task $b \Rightarrow a$
\task $a \Leftrightarrow b$
\task $\neg a \Leftrightarrow \neg b$
\end{tasks}
\end{exercise}

\begin{exercise}[Který výrok je tautologie?]
Určete, která z následujících výrokových formulí je tautologie:
\begin{tasks}(1)
\task $(A \land B) \Leftrightarrow (\neg A \lor \neg B)$
\task $(A \Leftrightarrow B) \Leftrightarrow (\neg A \Leftrightarrow \neg B)$
\task $(A \Rightarrow B) \Rightarrow (\neg A \Rightarrow \neg B)$
\end{tasks}
\end{exercise}

\begin{exercise}[Negace vět – přirozený jazyk]
Utuvořte negace následujících výroků:
\begin{tasks}(1)
\task „Číslo $72$ je dělitelné dvěma a třemi.“
\task „Mozart ani Beethoven nejsou čeští skladatelé.“
\task „\emph{Válku s mloky} napsal Karel Čapek nebo Alois Jirásek.“
\end{tasks}
\end{exercise}

\begin{exercise}[Negace implikací – matematické věty]
Utuvořte negace následujících výroků (pozor na tvar \,$\neg(P\Rightarrow Q)\Leftrightarrow (P\land\neg Q)$):
\begin{tasks}(1)
\task „Je-li $\sqrt{2}$ iracionální číslo, pak je iracionální i číslo $1+\sqrt{2}$.“
\task „Je-li trojúhelník rovnostranný, je rovnoramenný.“
\task „Není-li číslo $5$ přirozené, není přirozené ani číslo $1$.“
\end{tasks}
\end{exercise}

\begin{exercise}[Negace známých rčení]
Považujte následující rčení za výroky a zapište jejich negaci:
\begin{tasks}(1)
\task „Přišel jsem, viděl jsem, zvítězil jsem.“  
\task „Nebude-li pršet, nezmokneme.“  
\task „Já to platit nebudu, radši se dám na vojnu.“  
\task „Bude-li každý z nás z křemene, bude celý národ z kvádrů.“  
\end{tasks}
\end{exercise}

\begin{exercise}[Kvantifikátory ve větách]
Pomocí kvantifikátorů přepište následující tvrzení tak, aby vznikly pravdivé výroky:
\begin{tasks}(1)
\task „Pro čísla $x,y$ platí $x^2 + y^2 = 0$.“  
\task „Pro číslo $x$ platí $x^2 + 1 > 0$.“  
\end{tasks}
\end{exercise}


% =========================
% BONUS – ŠKOLNÍ KOLO (MO)
% =========================
\setcounter{section}{2}

\begin{exercise}[BONUS: Výroky o sourozencích]
Uvažujme 20 výroků dvou typů:

\begin{center}
\begin{tabular}{ll}
„Mám právě jednu sestru.“ & „Mám právě jednoho bratra.“ \\
„Mám právě dvě sestry.“   & „Mám právě dva bratry.“    \\
\vdots                    & \vdots                     \\
„Mám právě deset sester.“ & „Mám právě deset bratrů.“  \\
\end{tabular}
\end{center}

\begin{tasks}(1)
    \task Každý ze čtyř sourozenců pronesl jiný z těchto 20 výroků.  
          Je možné, že všichni čtyři měli pravdu?
    \task Najděte největší přirozené číslo $n$ takové, že každý z $n$ sourozenců může pronést jiný z těchto 20 výroků a mít pravdu.
\end{tasks}
\end{exercise}


% =========================
% BONUS
% =========================
\begin{exercise}[BONUS: Logická rekonstrukce trasy silnice]
Při stavbě silnice, která má spojovat dvě města, se rozhodovalo, 
kterými \textbf{obcemi} mezi nimi povede. Komise stanovila z 
ekonomických a technických důvodů tyto podmínky:
Silnice povede obcí $B$, povede-li obcí $A$.
Obcemi $A$ a $D$ povede, povede-li obcí $E$.
Silnice nepovede zároveň obcí $B$ a $C$.
Obcemi $C$ a $D$ buď povede oběma, anebo žádnou z nich.
Alespoň jednou z obcí $D$ a $E$ silnice povede.
    \begin{tasks}(1)
        \task Přepište podmínky do výrokové logiky.
        \task Rozhodněte, kterými obcemi silnice \emph{musí} a kterými \emph{nesmí} vést.
    \end{tasks}
\end{exercise}


% =========================
% ÚLOHA — Výroková logika (detektivní úloha)
% =========================
\begin{exercise}[BONUS: Logická úloha – Kdo je nevinný?]
    Někdo z podezřelých Antonín ($A$), Boris ($B$) a Cyril ($C$) v muzeu ukradl sošku.  
    Všichni tři svědci události vypověděli pravdivě následující výroky:

    \begin{enumerate}[label=\arabic*)]
        \item V muzeu nebyl podezřelý Cyril, nebo není pravda, že tam byl alespoň jeden z dvojice Antonín, Cyril.  
        \item Podezřelý Cyril byl v muzeu právě tehdy, když tam nebyl nikdo z dvojice Antonín, Boris.
        \item Pokud neplatí, že v muzeu byl Antonín spolu s Borisem, pak tam byl Cyril.
    \end{enumerate}

    Pomocí logických úvah (např. vytvořením pravdivostní tabulky nebo postupnou dedukcí) určete,  
    \textbf{kdo z podezřelých je nevinný}.
\end{exercise}


\end{document}
